% XeLaTeX can use any Mac OS X font. See the setromanfont command below.
% Input to XeLaTeX is full Unicode, so Unicode characters can be typed directly into the source.

% The next lines tell TeXShop to typeset with xelatex, and to open and save the source with Unicode encoding.

%!TEX TS-program = xelatex
%!TEX encoding = UTF-8 Unicode

\documentclass[a4paper,10pt]{report}
\usepackage[french]{babel}
\usepackage{standalone}
\usepackage{amsmath}
\usepackage{geometry}                % See geometry.pdf to learn the layout options. There are lots.
\geometry{a4paper,hmargin=2cm,vmargin=1.5cm,includeheadfoot}                   % ... or a4paper or a5paper or ... 
%\geometry{landscape}                % Activate for for rotated page geometry
%\usepackage[parfill]{parskip}    % Activate to begin paragraphs with an empty line rather than an indent
\usepackage{graphicx}
\usepackage{amssymb}
\usepackage{tkz-tab}
\usepackage{pgfplots}
\pgfplotsset{compat=1.18}
%\usepackage{xcolor}
%\definecolor{dauphineblue}{RGB}{47,68,134}


%\usepackage{fontspec,xltxtra,xunicode}
%\defaultfontfeatures{Mapping=tex-text}
%\setromanfont[Mapping=tex-text]{Hoefler Text}
%\setsansfont[Scale=MatchLowercase,Mapping=tex-text]{Gill Sans}
%\setmonofont[Scale=MatchLowercase]{Andale Mono}

\title{Math - Exercices}
\author{Poppy RAVEZ}
%\date{}                                           % Activate to display a given date or no date

\begin{document}
%!TEX TS-program = xelatex
%!TEX encoding = UTF-8 Unicode

\documentclass{report}
\usepackage[french]{babel}

%\usepackage{fontspec,xltxtra,xunicode}
%\defaultfontfeatures{Mapping=tex-text}
%\setromanfont[Mapping=tex-text]{Hoefler Text}
%\setsansfont[Scale=MatchLowercase,Mapping=tex-text]{Gill Sans}
%\setmonofont[Scale=MatchLowercase]{Andale Mono}

\usepackage{geometry}                % See geometry.pdf to learn the layout options. There are lots.
\geometry{a4paper,hmargin=2cm,vmargin=1.5cm}   
\usepackage{graphicx}
\usepackage{xcolor}
%\definecolor{dauphineblue}{cmyk}{100,85,15,0}
\definecolor{dauphineblue}{RGB}{47,68,134}
\author{Prénom NOM}

\begin{document}
\centering
\includegraphics[scale=0.08]{../../../logo/pdf/MIDO_centre}


\sffamily
\vspace{6cm}
\LARGE
\mdseries
Mathématiques

\huge
\bfseries
Raisonnement (3)

Ensembles

%\color{dauphineblue}
\vspace{3cm}
\rmfamily

\mdseries
\makeatletter
\large
Poppy RAVEZ

\vspace{0.2cm}
\normalsize
poppy.ravez@dauphine.psl.eu

\vspace{0.2cm}
Licence MIDO - 1\iere{} année

\vspace{5cm}
\@date
\makeatother


\end{document}
%%!TEX TS-program = xelatex
%!TEX encoding = UTF-8 Unicode

\documentclass{report}
\usepackage[french]{babel}

%\usepackage{fontspec,xltxtra,xunicode}
%\defaultfontfeatures{Mapping=tex-text}
%\setromanfont[Mapping=tex-text]{Hoefler Text}
%\setsansfont[Scale=MatchLowercase,Mapping=tex-text]{Gill Sans}
%\setmonofont[Scale=MatchLowercase]{Andale Mono}

\usepackage{geometry}                % See geometry.pdf to learn the layout options. There are lots.
\geometry{a4paper,hmargin=2cm,vmargin=1.5cm}   
\usepackage{graphicx}
\usepackage{xcolor}
%\definecolor{dauphineblue}{cmyk}{100,85,15,0}
\definecolor{dauphineblue}{RGB}{47,68,134}
\author{Prénom NOM}

\begin{document}
\centering
\includegraphics[scale=0.08]{../../../logo/pdf/MIDO_centre}


\sffamily
\vspace{6cm}
\LARGE
\mdseries
Mathématiques

\huge
\bfseries
Raisonnement (3)

Ensembles

%\color{dauphineblue}
\vspace{3cm}
\rmfamily

\mdseries
\makeatletter
\large
Poppy RAVEZ

\vspace{0.2cm}
\normalsize
poppy.ravez@dauphine.psl.eu

\vspace{0.2cm}
Licence MIDO - 1\iere{} année

\vspace{5cm}
\@date
\makeatother


\end{document}


\chapter*{Cours}

\subsection*{Application 2.3}

\subsubsection*{a. $\forall x \in \mathbb{R}, x \leq |x|$}

\begin{displaymath}
	\begin{cases}
		\forall x \in \mathbb{R}^{+}, |x| = x \geq x \\
		\forall x \in \mathbb{R}^{-}_{*}, |x| = -x > 0 > x
	\end{cases}
\end{displaymath}
Donc $\forall x \in \mathbb{R}, x \leq |x|$

\subsubsection*{b. $\forall x \in \mathbb{R}, |x| = 0 \Longleftrightarrow x = 0$}

Par définition : $x = 0 \Longrightarrow |x| = 0$.

Si $|x| = 0$, les solutions  vérifient :

\begin{displaymath}
	\begin{cases}
		x = |x| = 0 \text{ si } x \in \mathbb{R}^{+}\\
		-x = |x| = 0 \text{ si } x \in \mathbb{R}^{-}_{*}
	\end{cases}
\end{displaymath}

Donc $x=0$

L'implication réciproque $|x|=0, \Longrightarrow x=0$ est donc vraie ce qui démontre l'équivalence.

\subsubsection*{c. $\forall (x,y) \in \mathbb{R}^2, |xy| = |x|.|y|$}

\begin{tabular}{|c|c|c|}
	\hline
	$x$                  & $y$                  & $|xy|$\\
	\hline
	$\mathbb{R}^{+}$     & $\mathbb{R}^{+}$     & $|xy| = xy = |x|.|y|$ \\
	\hline
	$\mathbb{R}^{+}$     & $\mathbb{R}^{-}_{*}$ & $|xy| = -xy = x(-y) = |x|.|y|$ \\
	\hline
	$\mathbb{R}^{-}_{*}$ & $\mathbb{R}^{+}$     & $|xy| = -xy = (-x)y = |x|.|y|$ \\
	\hline
	$\mathbb{R}^{-}_{*}$ & $\mathbb{R}^{-}_{*}$ & $|xy| = xy =(-x)(-y) = |x|.|y|$ \\
	\hline
\end{tabular}

CQFD

Dans l'absolu et pour coller à la définition, il faudrait distinguer les cas $xy<0$ et $xy=0$ quand les deux variables n'ont pas
le même signe.


\subsubsection*{d. $\forall x \in \mathbb{R}, \forall z \in \mathbb{R}^{+}, |x| \leq z \Longleftrightarrow -z \leq x \leq z$}

Soit $z \in \mathbb{R}^{+}$ (et donc $-z \leq 0$).

Si $x \in \mathbb{R}^{+}$ (et donc $|x|=x \geq 0$):

\begin{itemize}
	\item $|x| \leq z \Longrightarrow -z \leq 0 \leq x \leq z \Longrightarrow -z \leq x \leq z$
	\item $-z \leq x \leq z \Longrightarrow -z \leq 0 \leq x \leq z \Longrightarrow |x| \leq z$
\end{itemize}

Ce qui démontre l'équivalence sur $\mathbb{R}^{+}$ ; l'équivalence sur $\mathbb{R}^{-}$ en est la conséquence
en remplaçant $x$ par $-x$ (les deux inégalités étant invariantes par changement de signe).

\subsection*{Application 2.12}

\subsubsection*{1. $\exp(-\frac{1}{x^2})$}
\begin{displaymath}
	\lim_{x \rightarrow 0} -\frac{1}{x^2} =-\infty \Longrightarrow 	\lim_{x \rightarrow 0} exp(-\frac{1}{x^2})= 0 \text{ car } \lim_{x \rightarrow -\infty} \exp(x) = 0
\end{displaymath}

\subsubsection*{2.}
\begin{displaymath}
	\lim_{x \rightarrow +\infty} (1 + e^x + x)= +\infty
\end{displaymath}

\subsubsection*{3.}
\begin{displaymath}
	\lim_{x \rightarrow +\infty} \frac{e^x+1}{e^x-1}= \lim_{x \rightarrow +\infty} \frac{1+e^{-x}}{1-e^{-x}} = 1 \text{ car }
	\lim_{x \rightarrow +\infty} e^{-x} = 0
\end{displaymath}


\subsection*{Application 2.16}

\begin{equation*}
	\begin{split}
		\ln(4-\sqrt{3}) + \ln(4+\sqrt{3})
		&= \ln ((4-\sqrt{3})(4+\sqrt{3}))\\
		&= \ln (4^2-\sqrt{3}^2) \\
		&= \ln 13
	\end{split}
\end{equation*}

\begin{equation*}
	\begin{split}
		\ln(\sqrt{x+1})
		&= 	\ln((x+1)^{1/2})\\
		&= \frac{\ln(x+1)}{2}
	\end{split}
\end{equation*}

\subsection*{Application 2.17}
D'après la définition de la dérivée d'une fonction :
\begin{displaymath}
\lim_{x \rightarrow 0} \frac{\ln(1+x)}{x} = \lim_{x \rightarrow 0} \frac{\ln(1+x) - \ln(1)}{x} = \ln'(1) = \frac{1}{1} = 1
\end{displaymath}


\subsection*{Application 2.24}

\subsubsection*{Application 2.24.i}

\begin{tikzpicture}
	\begin{axis}[
		axis lines=center,
		height=5cm,
		width=10cm,
		xlabel={$x$},
		ylabel={$f(x)$},
		xmin=0,
		ymin=0, 
		]
		\addplot[blue,domain=0:8,samples=161]{x^(1/3))};
	\end{axis}	
\end{tikzpicture}

Comme vu dans la proposition 2.22, $p_{1/3}$ est étendue à $\mathbb{R}$ comme la réciproque de $p_3$ qui est une bijection
 de  $\mathbb{R}$ dans $\mathbb{R}$ : la racine cubique. On peut formellement définir $p_{1/3}$ sur $\mathbb{R}^-$ par $\forall x \in \mathbb{R}^-,
 p_{1/3}(x) = -p_{1/3}(-x)$. On complète donc le graphe par son symétrique par rapport à l'origine.

\begin{tikzpicture}
	\begin{axis}[
		axis lines=center,
		height=5cm,
		width=10cm,
		xlabel={$x$},
		ylabel={$f(x)$},
		]
		\addplot[blue,domain=-8:0,samples=161]{-(-x)^(1/3))};
		\addplot[blue,domain=0:8,samples=161]{x^(1/3))};
	\end{axis}	
\end{tikzpicture}

\subsubsection*{Application 2.24.ii}

$D=\mathbb{R}^{+}_*$

\begin{tikzpicture}
	\begin{axis}[
		axis lines=center,
		height=5cm,
		width=10cm,
		xlabel={$x$},
		ylabel={$f(x)$},
		xmin=0,
		ymin=0, 
		]
		\addplot[blue,domain=0.1:16,samples=161]{x^(-1/4))};
	\end{axis}	
\end{tikzpicture}

\subsection*{Application 2.26 - Pour aller plus loin}
Cela n'est vrai sur $\mathbb{R}$ que pour les entiers impairs car $x^{2k}$ n'est pas une bijection.

\subsection*{Application 2.32}
\begin{equation*}
	\begin{split}
		\forall x \in \mathbb{R}^{+}_{*}, \left(  x^{\frac{1}{x}} \right)^{x^2+2x}
		&= x^{\frac{x^2+2x}{x}} \\
		&= x^{x+2}
	\end{split}
\end{equation*}


\subsection*{Application 2.34 - Dérivabilité de $p_{\alpha}$ en $0$ avec $\alpha \geq 1$}

$p_{\alpha}$ est définie sur $\mathbb{R}^{+}_{*}$ par $p_{\alpha}(x) = \exp(\alpha \ln(x))$, cependant :
\begin{displaymath}
	\lim_{x\rightarrow 0^{+}} \ln(x) = -\infty \Longrightarrow\forall \alpha>0, \lim_{x\rightarrow 0^{+}} p_{\alpha}(x) = 0
\end{displaymath}

On peut donc, si $\alpha > 0$, définir par continuité la fonction $p_{\alpha}$ en $0$ avec la valeur nulle.

On a alors :
\begin{equation*}
	\begin{split}
		\lim_{h\rightarrow 0^{+}} \frac{p_{\alpha}(h) - 0}{h}
		&= \lim_{h\rightarrow 0^{+}} \frac{e^{\alpha \ln(h)}}{h} \\
		&= \lim_{h\rightarrow 0^{+}} \frac{e^{\alpha \ln(h)}}{e^{\ln(h)}} \\
		&= \lim_{h\rightarrow 0^{+}} e^{(\alpha-1) \ln(h)} \\
		&=0 \text{ si $\alpha >1$ car alors $\lim_{h\rightarrow 0^{+}} (\alpha-1) \ln(h) = -\infty$}
	\end{split}
\end{equation*}

La fonction $p_{\alpha}$ prolongée en $0$ avec $p_{\alpha}(0)=0$ est donc dérivable à droite en $0$ si $\alpha>1$
avec $p'_{\alpha}(0)=0$.


\subsection*{Application 2.41 - Limite en $0$ de $x^x$}

\begin{equation*}
	\begin{split}
		\lim_{x\rightarrow 0^{+}} x^x
		&= \lim_{x\rightarrow 0^{+}} e^{x\ln(x)} \\
	    &= 1 \text{ car  $\lim_{x\rightarrow 0^{+}} x\ln(x) = 0$ d'après la proposition 2.39}
	\end{split}
\end{equation*}

\chapter*{Exercices}

\subsection*{Exercice 2.1}

\begin{equation*}
	\begin{split}
		\forall x \in \mathbb{R}, \lim_{h\rightarrow 0} \dfrac{\dfrac{1}{1+(x+h)^2}-\dfrac{1}{1+x^2}}{h}
		&= \lim_{h\rightarrow 0} \dfrac{1+x^2-1-(x+h)^2}{h(1+(x+h)^2)(1+x^2)} \\
		&= \lim_{h\rightarrow 0} \dfrac{-2xh-h^2}{h(1+(x+h)^2)(1+x^2)} \\
		&= \lim_{h\rightarrow 0} \dfrac{-2x-h}{(1+(x+h)^2)(1+x^2)} \\
		&= \frac{-2x}{(1+x^2)^2} \\
	\end{split}
\end{equation*}

Car la limite du dénominateur est un réel strictement positif.

La fonction $f$ est donc dérivable sur $\mathbb{R}$ et $f'(x) = \frac{-2x}{(1+x^2)^2}$

On constate de manière évidente que $f$ est paire, strictement positive, atteint son maximum en $0$ et que :
\begin{displaymath}
	\lim_{x \rightarrow \pm \infty} f(x) = 0^{+}
\end{displaymath}


On en déduit le tableau de variation suivant :

\begin{tikzpicture}
	\tkzTabInit{$x$ / 1 , $f'(x)$ / 1, $f(x)$/2}{$-\infty$ ,$0$, $+\infty$}
	\tkzTabLine{, + , 0 , - , }
	\tkzTabVar{-/$0$, +/$1$, -/$0$}
\end{tikzpicture}

Recherchons les points d'inflexion (où elle traverse sa tangente) caractérisés par $f''(x)=0$ :

\begin{equation*}
	\begin{split}
		f''(x) &= \dfrac{-2(1+x^2)^2-2(-2x)(2x)(1+x^2)}{(1+x^2)^4} \\
		       &= \dfrac{-2(1+x^2)+8x^2}{(1+x^2)^3} \\
		       &= 2\dfrac{3x^2-1}{(1+x^2)^3} 
	\end{split}
\end{equation*}

Les points d'inflexion sont donc $\binom{\pm\frac{1}{\sqrt{3}}}{\frac{3}{4}}$.

\begin{tikzpicture}
	\begin{axis}[
		axis lines=center,
		height=5cm,
		width=10cm,
		xlabel={$x$},
		ylabel={$f(x)$},
		xmin=-4, xmax= 4,
		ymin=0, ymax=1.1,
		]
		\addplot[blue,samples=81]{1/(1+x^2)};
	\end{axis}	
\end{tikzpicture}


\subsection*{Exercice 2.2}

$f$ est impaire, continue et dérivable sur $\mathbb{R}$.

$f(0) = 0$.

\begin{displaymath}
\lim_{x \rightarrow -\infty} f(x) = 0^{-}
\end{displaymath}

\begin{displaymath}
	\lim_{x \rightarrow +\infty} f(x) = 0^{+}
\end{displaymath}

On peut remarquer que :
\begin{displaymath}
	f_n\left(\frac{1}{\sqrt{n}} x \right) = \frac{1}{\sqrt{n}} f_1(x).
\end{displaymath}

La représentation graphique de $f_n$ est donc la transformée de celle de $f_1$
par l'homothétie ayant pour centre l'origine et de rapport $\frac{1}{\sqrt{n}}$.
L'homothétie conservant les directions (et donc les sens de variations, extrema et
points d'inflexion, et pentes des tangentes), ll suffit d'étudier $f_1$.

Cette caractéristique, si elle nous échappe initialement, deviendra évidente durant
l'étude de la fonction lorsque l'on constatera l'alignement des points particuliers
(extrema et points d'inflexion).

\subsubsection*{Variations}

\begin{displaymath}
	f'(x) = \frac{1+nx^2-2nxx}{(1+nx^2)^2}=\frac{1-nx^2}{(1+nx^2)^2}
\end{displaymath}

$f'(0) = 1$ : toutes les courbes en la même tangente à l'origine

Les extrema (où $f'(x)=0$) sont donc $\pm\frac{1}{\sqrt{n}}$

$f(\frac{1}{\sqrt{n}}) = \frac{1}{2\sqrt{n}}$

On constate que $\frac{f(1/\sqrt{n})}{1/\sqrt{n}} = \frac{1}{2}$, les extrema sont donc tous sur la droite
d'équation $y=\frac{x}{2}$.

Le tableau de variation est donc :

\begin{tikzpicture}
	\tkzTabInit{$x$ / 1 , $f'(x)$ / 1, $f(x)$/2}{$-\infty$ ,$-\frac{1}{\sqrt(n)}$, $+\frac{1}{\sqrt(n)}$, $+\infty$}
	\tkzTabLine{, - , 0 , +, 0, - , }
	\tkzTabVar{+/$0^{-}$, -/$-\frac{1}{2\sqrt{n}}$, +/$\frac{1}{2\sqrt{n}}$, -/$0^{+}$}
	\tkzTabVal{2}{3}{0.5}{0}{0}
\end{tikzpicture}

\subsubsection*{Points d'inflexion}

\begin{displaymath}
	f''(x) = \frac{-2nx(1+nx^2)^2-4nx(1-nx^2)(1+nx^2)}{(1+nx^2)^4}
	       =-2nx\frac{1+nx^2+2(1-nx^2)}{(1+nx^2)^3}
	       = -2nx\frac{3-nx^2}{(1+nx^2)^3}
\end{displaymath}

Chaque courbe présente donc 3 points d'inflexion alignés en $-\sqrt{\frac{3}{n}}$, $0$ et $\sqrt{\frac{3}{n}}$.

On constate que $\frac{f(\sqrt{\frac{3}{n}})}{\sqrt{\frac{3}{n}}} = \frac{1}{4}$ et que $f'\left({\pm \sqrt{\frac{3}{n}}}\right)=-1/8$, ce qui montre que tous les points d'inflexion sont sur la droite
$y=x/4$ et que les tangentes en ces points sont toutes parallèles.

\subsubsection*{Représentation graphique}

\begin{tikzpicture}
	\begin{axis}[
		axis lines=center,
		height=5cm,
		width=10cm,
		xlabel={$x$},
		ylabel={$f(x)$},
		xmin=-4, xmax= 4,
		ymin=-0.6, ymax=0.6,
		legend pos= north west,
		]
		\addplot[blue,samples=81]{x/(1+x^2)};
		\addlegendentry{$n=1$}
		\addplot[red,samples=81]{x/(1+2*x^2)};
		\addlegendentry{$n=2$}
		\addplot[green,samples=81]{x/(1+3*x^2)};
		\addlegendentry{$n=3$}
	\end{axis}	
\end{tikzpicture}

\subsection*{Exercice 2.3}

\subsubsection*{1.}

\begin{equation*}
	\begin{split}
		|\sin^{12}(x) - 2\cos^9(x)| &\leq |\sin^{12}(x)|+ |2\cos^9(x)| \text{ (inégalité triangulaire)} \\
		                         &\leq 1 + 2 \\
		                         & \leq 3
	\end{split}
\end{equation*}

\subsubsection*{2.}

\begin{equation*}
	\begin{split}
		|1-x^{12} +4x^3| &\leq |1- x^{12}|+ |4x^3| \text{ (inégalité triangulaire)} \\
		&\leq 1 + 4 \text{ car }x \in [-1, 1], 1- x^{12} \in [0 ; 1]\\
		& \leq 5
	\end{split}
\end{equation*}


\subsection*{Exercice 2.4}

\subsubsection*{(a)}

Les deux termes étant positifs, on peut élever l'égalité au carré ou raisonner par disjonction de cas en fonction du signe de $x+1$ :

\begin{alignat*}{2}
	&                    & |x+1| &=2 \\
	&\Longleftrightarrow &(x+1)^2 &= 4 \\
	&\Longleftrightarrow & x^2 +2x -3 &= 0 \\
	&\Longleftrightarrow & (x-1)(x+3) &= 0 \\
\end{alignat*}
\begin{alignat*}{2}
&                    & |x+1| &=2 \\
&\Longleftrightarrow &(x+1=2 \wedge x+1 \geq 0) &\vee (-x-1=2 \wedge x+1 \leq 0) \\
&\Longleftrightarrow &(x=1 \wedge x \geq -1) &\vee (x=-3 \wedge x \leq -1) \\
&\Longleftrightarrow & x=1 &\vee x=-3 \\
\end{alignat*}

Donc : $S=\left\lbrace -3 ; 1 \right\rbrace$

\subsubsection*{(b)}


Les deux termes étant positifs, on peut élever au carré :

\begin{alignat*}{2}
	&                    & |x+3| &\leq 4 \\
	&\Longleftrightarrow &(x+3)^2 &\leq 16 \\
	&\Longleftrightarrow & (x^2 + 6x -7) & \leq 0 \\
	&\Longleftrightarrow & (x+7)(x-1) \leq 0\\
\end{alignat*}

L'ensembe des solutions est donc l'intervalle fermé entre les 2 racines $-7$ et $1$.

Ou par disjonction de cas :

\begin{alignat*}{2}
	&                    & |x+3| &\leq 4 \\
	&\Longleftrightarrow &(x+3\leq 4 \wedge x+3 \geq 0) &\vee (-x-3\leq 4 \wedge x+3 \leq 0) \\
	&\Longleftrightarrow & (-3 \leq x \leq 1) &\vee (-7 \leq x \leq -3) \\
	&\Longleftrightarrow & -7 \leq &x \leq 1 \\
\end{alignat*}

Donc : $S=[-7 ; 1]$

\subsubsection*{(c)}


Les deux termes étant positifs, on peut élever au carré :

\begin{alignat*}{2}
	&                    & |x+1| &> 9 \\
	&\Longleftrightarrow &(x+1)^2 &> 9 \\
	&\Longleftrightarrow & x^2 + 2x -8 & >0 \\
	&\Longleftrightarrow & (x+4)(x-2) &> 0\\
\end{alignat*}

L'ensembe des solutions est donc $\mathbb{R}$ diminué de l'intervalle fermé entre les 2 racines $-4$ et $2$ : $\mathbb{R} \setminus [-4 ; 2]$

Ou par disjonction de cas :

\begin{alignat*}{2}
	&                    & |x+1| &>3  \\
	&\Longleftrightarrow &(x+1 > 3 \wedge x+1 \geq 0) &\vee (-x-1 >3 \wedge x+1 \leq 0) \\
	&\Longleftrightarrow & (x > 2) &\vee (x < -4)
\end{alignat*}

Donc : $S=]-\infty ; -4] \cup ]2 ; +\infty[$

\subsubsection*{(d)}


Les deux termes étant positifs, on peut élever au carré :

\begin{alignat*}{2}
	&                    & |2x-4| &= |x+2| \\
	&\Longleftrightarrow &(2x-4)^2 &= (x+2)^2 \\
	&\Longleftrightarrow & 3x^2 -20x +12 & =0 \\
	&\Longleftrightarrow & (3x-2)(x-6) &= 0\\
\end{alignat*}


Ou par disjonction de cas :

\begin{tikzpicture}
	\tkzTabInit{$x$ / 1   , $|2x-4|$ / 1, $|x+2|$/1, Equation/1, S/1}
	           {$-\infty$ ,$-2$, $2$    , $+\infty$}
	\tkzTabLine{, -2x+4       , , -2x+4, , 2x-4 , }
	\tkzTabLine{, -x-2        , , x+2, , x+2 , }
	\tkzTabLine{, -x+6=0      , , -3x+2=0, , x-6=0 , }
	\tkzTabLine{, \varnothing ,\cup , \left\lbrace \frac{2}{3}\right\rbrace ,\cup , \left\lbrace 6 \right\rbrace , }
\end{tikzpicture}

Donc : $S=\left\lbrace \frac{2}{3} ; 6 \right\rbrace$

\subsubsection*{(e)}

Les deux termes étant positifs, on peut élever au carré :

\begin{alignat*}{2}
	&                    & |2x+4|        &\leq |x+1| \\
	&\Longleftrightarrow & (2x+4)^2      &\leq (x+1)^2 \\
	&\Longleftrightarrow & 3x^2 +14x +15 &\leq 0 \\
	&\Longleftrightarrow & (3x+5)(x+3)   &\leq 0 \\
	&\Longleftrightarrow & -3 \leq &x \leq -\frac{5}{3}
\end{alignat*}

Ou par disjonction de cas :

\begin{tikzpicture}
	\tkzTabInit{$x$ / 1   , $|2x+4|$ / 1, $|x+1|$/1,Equation/1, S/1}
	{$-\infty$ ,$-2$, $-1$    , $+\infty$}
	\tkzTabLine{, -2x-4 , , 2x+4, , 2x+4 , }
	\tkzTabLine{, -x-1 , , -x-1, , x+1 , }
	\tkzTabLine{, -x-3 \leq 0 , , 3x+5 \leq 0, , x+3 \leq 0 , }
	\tkzTabLine{, [-3 ;-2] ,\cup , [-2 ; -\frac{5}{3}] ,\cup , \varnothing , }
\end{tikzpicture}

Donc : $S=\left[ -3 ; -\frac{5}{3} \right] $

\subsubsection*{(f)}

Si $|x|<1$, l'égalité ne peut être vérifiée car le terme de droite est strictement négatif et celui de
gauche positif. Considérons donc que $|x|\geq 1$ et élevons au carré :

\begin{alignat*}{2}
	&\forall x \in \mathbb{R} \setminus ]-1;1[,  & |x+1|   &= |x| -1 \\
	&\Longleftrightarrow                         & (x+1)^2 &= (|x|-1)^2 \\
	&\Longleftrightarrow                         & x       &= -|x| \\
	&\Longleftrightarrow                         & x       &\leq 0 \\
		&\Longleftrightarrow                     & x       &\leq -1 \\
\end{alignat*}



Ou par disjonction de cas :

\begin{tikzpicture}
	\tkzTabInit{$x$ / 1   , $|x+1|$ / 1, $|x|$/1,Equation/1, S/1}
	{$-\infty$ ,$-1$, $0$    , $+\infty$}
	\tkzTabLine{, -x-1 , , x+1, , x+1 , }
	\tkzTabLine{, -x , , -x, , x , }
	\tkzTabLine{, 0=0 , , 2x+2=0, , 2=0 , }
	\tkzTabLine{, ]-\infty ;-1] ,\cup , \left\lbrace -1 \right\rbrace , \cup, \varnothing , }
\end{tikzpicture}

Donc : $S= ]-\infty ; -1]$

\subsection*{Exercice 2.5}


Par disjonction de cas :

\begin{tikzpicture}
	\tkzTabInit{$x$ / 1   , $|2x-1|$ / 1, $|3x+1|$/1,Equation/1, S/1}
	           {$-\infty$ ,$-\frac{1}{3}$, $\frac{1}{2}$    , $+\infty$}
	\tkzTabLine{, -2x+1  , , -2x+1  , , 2x-1    , }
	\tkzTabLine{, -3x-1  , ,  3x+1  , , 3x+1    , }
	\tkzTabLine{, 4x+3=1 , , -8x-1=1, , -4x-3=1 , }
	\tkzTabLine{,   \left\lbrace -\frac{1}{2} \right\rbrace,\cup , \left\lbrace -\frac{1}{4} \right\rbrace , \cup, \varnothing , }
\end{tikzpicture}

Donc : $S=\left\lbrace -\frac{1}{2} ; -\frac{1}{4} \right\rbrace$

\subsection*{Exercice 2.6}

\begin{equation*}
	\begin{split}
		\forall x \in [-3 ; 1],
		& \begin{cases}
			0 \leq |x+1| \leq 2 \\
		    -1 \leq x+2 \leq 3 \\
		    1 \leq x^2 + 1 \leq 10
		\end{cases} \\
		\Longrightarrow 
		& \begin{cases}
			-2 \leq |x+1|(x+2) \leq 6 \\
			0 < \frac{1}{x^2 + 1} \leq 1
		\end{cases} \\
		\Longrightarrow 
		& -2 \leq \frac{|x+1|(x+2)}{x^2 + 1} \leq 6 
	\end{split}
\end{equation*}

\subsection*{Exercice 2.7}

\subsubsection*{Rappel de logique}

«$A$ est vraie si et seulement si (ssi) $B$ est vraie» signifie que A et B sont équivalentes ($A\Longleftrightarrow B$).

«$A$ vraie si $B$ vraie » signifie que $B \Rightarrow A$

«$A$ vraie seulement si $B$ vraie » signifie que  $\bar{B} \Rightarrow \bar{A}$ soit $A \Rightarrow B$

On démontre une équivalence soit par équivalences successives soit en montrant que $A\Rightarrow B$ et que 
$B\Rightarrow A$.

Une implication est équivalente à sa contraposée : on peut démontrer que $A\Rightarrow B$ en montrant que
$\bar{B}\Rightarrow \bar{A}$ ou $\bar{A}$ est la négation de $A$.

\subsubsection*{1. Egalité de l'inégalité triangulaire}


L'égalité est évidemment vraie si $x$ et $y$ sont de même signe :
\begin{displaymath}
	xy \geq 0 \Longrightarrow |x+y| = |x| + |y|
\end{displaymath}

Plutôt que de prouver la réciproque ($|x+y| = |x| + |y| \Longrightarrow xy \geq 0$), démontrons
la contraposée : $xy < 0 \Longrightarrow |x+y| \neq |x| + |y|$.

Soient $x$ et $y$ tels que $ x < 0 < y$, on alors :

\begin{displaymath}
		|x+y| - |x| - |y|
		= \begin{cases}
			x+y-(-x)-y = 2x \neq 0 \text { si } y \geq -x \\
			-x-y-(-x)-y = -2y \neq 0 \text { si } y \leq -x \\
		\end{cases}
\end{displaymath}

L'égalité n'est donc pas vérifiée si les 2 variables n'ont pas le même signe ce qui démontre la contraposée
et donc l'équivalence.

\subsubsection*{2. min et max}

Les 2 fonctions étant symétriques en $x$ et $y$ (car $|x-y| = |y-x]$), il suffit de vérifier les égalités
pour $x \leq y$ et elles seront vraies pour $x \geq y$.

Supposons $x \leq y$ (et donc $|x-y|=y-x$), alors :
\begin{equation*}
	\begin{split}
		\frac{x + y + |x-y|}{2} &= \frac{x + y + (y-x)}{2} \\
		                        &= y \\
		                        &= max(x,y)
	\end{split}
\end{equation*}

et

\begin{equation*}
	\begin{split}
		\frac{x + y - |x-y|}{2} &= \frac{x + y - (y-x)}{2} \\
		&= x \\
		&= min(x,y)
	\end{split}
\end{equation*}

\subsection*{Exercice 2.8}

\subsubsection*{1.}
\begin{equation*}
	\begin{split}
		&\frac{\exp(x^2-1)}{\exp(x+1)} \leq 1 \\
		\Longleftrightarrow &\exp(x^2-1 -(x+1)) \leq 1 \\
		\Longleftrightarrow &\exp((x+1)(x-2)) \leq 1 \\
		\Longleftrightarrow &(x+1)(x-2) \leq 0 \\
		\Longleftrightarrow & -1 \leq x \leq 2 \\
	\end{split}
\end{equation*}

Donc : $S= [-1 ; 2]$

\subsubsection*{2.}
\begin{equation*}
	\begin{split}
		& e^{2x} + e^x -6 = 0 \\
		\Longleftrightarrow & e^x = \frac{-1\pm \sqrt{25}}{2} \\
		\Longrightarrow & e^x = 2 \text{ car $e^x$ est strictement positif} \\
		\Longrightarrow & x = \ln 2
	\end{split}
\end{equation*}

Donc : $S=\left\lbrace \ln 2 \right\rbrace$

\subsubsection*{3.}
\begin{equation*}
	\begin{split}
		\forall x \in \mathbb{R} \setminus \left\lbrace -1 ; 1 \right\rbrace, & \ln |x+1| - \ln |x-1| \leq \ln 2 \\
		\Longleftrightarrow & \ln \left| \frac{x+1}{x-1} \right| \leq \ln 2 \\
		\Longleftrightarrow & \left| \frac{x+1}{x-1} \right| \leq 2 \text{ car $\ln$ est croissante sur $\mathbb{R}^{+}_{*}$} \\
		\Longleftrightarrow & \left( \frac{x+1}{x-1} \right)^2 \leq 4 \text{ car les 2 termes sont positifs} \\
		\Longleftrightarrow & 3x^2-10x+3 \geq 0 \\
		\Longleftrightarrow & (3x-1)(x-3) \geq 0 \\
		\Longleftrightarrow & x \in ]-\infty ; \frac{1}{3}] \cup [3 : +\infty[ \\
	\end{split}
\end{equation*}

Et donc $S = ]-\infty ; -1[ \cup ]-1 ; \frac{1}{3}] \cup [3 ; +\infty[$


On aurait également pu étudier la fonction $f(x) = \left| \frac{x+1}{x-1} \right|$ définie sur $\mathbb{R} \setminus \left\lbrace -1 ; 1 \right\rbrace$.

\begin{displaymath}
	\lim_{x \rightarrow -\infty} f(x) = 1^{-}
\end{displaymath}

\begin{displaymath}
	\lim_{x \rightarrow -1} f(x) = 0
\end{displaymath}

\begin{displaymath}
	\lim_{x \rightarrow 1-} f(x) = +\infty
\end{displaymath}

\begin{displaymath}
	\lim_{x \rightarrow 1+} f(x) = +\infty
\end{displaymath}

\begin{displaymath}
	\lim_{x \rightarrow +\infty} f(x) = 1^{+}
\end{displaymath}

\begin{tikzpicture}
	\tkzTabInit{$x$ / 1   , $|x+1|$ / 1, $|x-1|$/1,$f(x)$/1, $f'(x)$/1, Variations/2}
	           {$-\infty$  ,$-1$ , $1$    , $+\infty$}
	\tkzTabLine{, -x-1               , d , x+1               , d , x+1    , }
	\tkzTabLine{, -x+1               , d , -x+1              , d , x-1    , }
	\tkzTabLine{, \frac{x+1}{x-1}    , d , -\frac{x+1}{x-1}  , d , \frac{x+1}{x-1} , }
	\tkzTabLine{, -\frac{2}{(x-1)^2} , d , \frac{2}{(x-1)^2} , d , -\frac{2}{(x-1)^2} , }
	\tkzTabVar{+/$1^{-}$, -C/$0$ , +D+/$+\infty$/$+\infty$, -/$1^{+}$ , }
\end{tikzpicture}

L'équation $f(x)=2$ présente donc 1 racine sur l'intervalle $]-1 ; 1[$ et une autre sur $]1; +\infty[$ : $\frac{1}{3}$ et $3$.

On a donc $S = ]-\infty ; -1[ \cup ]-1 ; \frac{1}{3}] \cup [3 ; +\infty[$

\subsection*{Exercice 2.9}

Soit $f(x) = \frac{xe^x}{x+1}$

Il s'agit du produit de fonctions ($x$, $e^x$ et $1/(x+1)$) dérivables sur $\mathbb{R}\setminus\{-1\}$,
elle est donc dérivable sur ce même domaine. 

\begin{equation*}
	\begin{split}
		f'(x) &= \frac{(e^x + x e^x)(x+1)-x e^x}{(x+1)^2} \\
		      &= \frac{(x+1)^2-x}{(x+1)^2}e^x \\
		      &= \frac{x^2 + x + 1}{(x+1)^2}e^x
	\end{split}
\end{equation*}

\subsection*{Exercice 2.10}

Soit $f(x) = e^\frac{^1}{\ln x}$

$\exp$ étant défini sur $\mathbb{R}$ et $\ln$ sur $\mathbb{R}^{+}_{*}$ mais ne devant pas être nul (car au dénominateur d'une
fraction), on a $D = \mathbb{R}^{+}_{*}\setminus\{1\}$.

\begin{displaymath}
	\lim_{x \rightarrow 0^{+}} f(x) = 1^{-}
\end{displaymath}

\begin{displaymath}
	\lim_{x \rightarrow 1^{-}} f(x) = 0
\end{displaymath}

\begin{displaymath}
	\lim_{x \rightarrow 1^{+}} f(x) = +\infty
\end{displaymath}

\begin{displaymath}
	\lim_{x \rightarrow +\infty} f(x) = 1^{+}
\end{displaymath}

\begin{displaymath}
	f'(x) =-\frac{1}{x\ln^2 x} e^{\frac{1}{\ln x}} < 0
\end{displaymath}

$f$ est donc strictement décroissante sur $D$.


On peut la prolonger en $0$ par continuité avec $f(0) = 1$

\begin{tikzpicture}
	\tkzTabInit{$x$/1, $f'$/1, $f$/3 }
	           {$0$  , $1$   , $+\infty$}
	\tkzTabLine{d, - , d , -, }

	\tkzTabVar{+C/$1^{-}$, -D+/$0$/+$\infty$ , -/$1$ }
\end{tikzpicture}

\begin{tikzpicture}
	\begin{axis}[
		axis lines=center,
		height=6cm,
		width=10cm,
		xlabel={$x$},
		ylabel={$f(x)$},
		ymin=0,
		]
		\addplot[blue, domain=0:0.9, samples=81]{exp(1/ln(x))};
		\addplot[blue, domain=1.5:5, samples=51]{exp(1/ln(x))};
	\end{axis}	
\end{tikzpicture}

\subsection*{Exercice 2.11}
\subsubsection*{1. Montrons que $x \geq \ln(1+x)$ sur $\mathbb{R}^+$}

Soit $f(x) = x - \ln(1+x)$

$f(0) = 0$

\begin{displaymath}
	f'(x) = 1 - \frac{1}{1+x} = \frac{x}{1+x} \geq 0 \text{ sur } \mathbb{R}^{+}
\end{displaymath}

$f$ est donc croissante et par conséquence positive sur $\mathbb{R}^{+}$ : $x \geq \ln(1+x)$


\subsubsection*{2. Montrons que $x -\frac{x^2}{2} \leq \ln(1+x)$ sur $\mathbb{R}^+$}

Soit $g(x) = \ln(1+x) + \frac{x^2}{2} - x$

$g(0) = 0$
\begin{displaymath}
	g'(x) = \frac{1}{1+x} + x - 1 = \frac{x^2}{1+x} \geq 0 \text{ sur } \mathbb{R}^{+}
\end{displaymath}

$g$ est donc croissante et positive sur $\mathbb{R}^{+}$ soit : $\ln(1+x) \geq -\frac{x^2}{2} + x$

\subsection*{Exercice 2.12}

$f(x) = \frac{1}{\sqrt{x^2 +x - 6}}$

Le domaine de définition est le domaine où $x^2 + x -6$ est strictement positive soit $\mathbb{R}$ diminué de l'éventuel
intervalle fermé entre ses racines ($-3$ et $2$).

$D = ]-\infty ; -3 [ \cup ]2; +\infty[$

\subsection*{Exercice 2.13}
\subsubsection*{1. Dérivabilité de $p_\alpha$ en 0}

Soit $\alpha \in ]0;1[$

\begin{equation*}
	\begin{split}
		\lim_{x \rightarrow 0^{+}} \frac{p_{\alpha}(x) - p_{\alpha}(0)}{x}
		&= \lim_{x \rightarrow 0^{+}} \frac{e^{\alpha \ln(x)}}{e^{\ln x}} \\
		&= \lim_{x \rightarrow 0^{+}} e^{(\alpha-1) \ln(x)} \\
		&= +\infty
	\end{split}
\end{equation*}

$p_\alpha$ n'est donc pas dérivable en $0$.

\subsubsection*{2.}

C'est la conséquence du cas précédent $\alpha = \frac{1}{2}$.

On peut calculer :

\begin{equation*}
	\begin{split}
		\lim_{h \rightarrow 0^{+}} \frac{f(1-h) - f(1)}{-h}
			&= \lim_{h \rightarrow 0^{+}} -\frac{\sqrt{h(2-h)}}{h} \\
			&= \lim_{h \rightarrow 0^{+}} -\sqrt{\frac{(2-h)}{h}} \\
			&= -\infty
	\end{split}
\end{equation*}

$f$ n'est donc pas dérivable en $1$.

Le calcul est identique en $-1$ mais $f$ étant paire la non dérivabilité en $-1$ est évidente.


\subsection*{Exercice 2.14}

Soit $f(x) = \sqrt{1-x^2}$

Domaine de définition : $D = [-1 ; 1]$

$f'(x) = -\frac{x}{\sqrt{1-x^2}}$

\begin{tikzpicture}
	\tkzTabInit{$x$/1, $f'$/1, $f$/2 }
	{$-1$  , $0$   , $1$}
	\tkzTabLine{, + , 0 , -, }
	
	\tkzTabVar{-/$0$, +/$1$ , -/$0$ }
\end{tikzpicture}

\begin{tikzpicture}
	\begin{axis}[
		axis lines=center,
		height=5cm,
		width=10cm,
		xlabel={$x$},
		ylabel={$f(x)$},
		xmin=-1.1, xmax= 1.1,
		ymin=0, ymax=1.1,
		]
		\addplot[blue,samples=161]{sqrt(1-x^2)};
	\end{axis}	
\end{tikzpicture}

On a : $x^2 + f^2(x) = 1$ ; la courbe est donc le demi-cercle centré sur l'origine et de rayon 1.

\subsection*{Exercice 2.15}

Soient $f(x) = \sqrt{x}$ et $g(x)= \ln x$.

\begin{displaymath}
	D_{f \circ g} = [1 ; +\infty[
\end{displaymath}

\begin{displaymath}
	D_{g \circ f} = \mathbb{R}^{+}_{*}
\end{displaymath}


\subsection*{Exercice 2.16}
\subsubsection*{1.}
$f(x) = \sqrt{x-1} - \sqrt{2x-3}$ est continue sur son domaine de définition.

$D = [\frac{3}{2} ; +\infty[$

$f(\frac{3}{2}) = \frac{1}{\sqrt{2}}$

$f(x) = 0 \Rightarrow x-1 = 2x-3 \Rightarrow x = 2 $

$\lim_{x \rightarrow +\infty} f(x) = -\infty$

On en déduit le tableau suivant :

\begin{tikzpicture}
	\tkzTabInit{$x$/1, $f$/1 }
	{$\frac{3}{2}$  , $2$   , $+\infty$}
	\tkzTabLine{, + , z , -, }
\end{tikzpicture}

\subsubsection*{2.}
$f(x) = \sqrt{|x-1|} - \sqrt{|2x-3|}$

$D = \mathbb{R}$

$f(x) = 0 \Longleftrightarrow (x-1)^2 = (2x-3)^2 \Longleftrightarrow 3x^2 -10x +8 =0 \Longleftrightarrow (x-2)(3x-4) =0$

$f(\frac{3}{2}) = \frac{1}{\sqrt{2}} > 0$

$\lim_{x \rightarrow -\infty} f(x) = \lim_{x \rightarrow +\infty} f(x) = -\infty$

\begin{tikzpicture}
	\tkzTabInit{$x$/1, $f$/1 }
	{$-\infty$, $\frac{4}{3}$  , $2$   , $+\infty$}
	\tkzTabLine{, - , z , +, z, - }
\end{tikzpicture}


\subsubsection*{3.}
\begin{equation*}
	\begin{split}
		f(x) &= (x^2-3)(2-\sqrt{x})(\sqrt{x^2-2x+1}-4) \\
		     &= (x^2-3)(2-\sqrt{x})(\sqrt{(x-1)^2}-4) \\
		     &= (x^2-3)(2-\sqrt{x})(|x-1|-4) \\
	\end{split}
\end{equation*}

$D = \mathbb{R}^{+}$

$f(x) = 0 \Longleftrightarrow x \in \{ \sqrt{3} ; 4 ; 5 \}$

$\lim_{x \rightarrow +\infty} f(x) = -\infty$

\begin{tikzpicture}
	\tkzTabInit{$x$/1, $x^2-3$/1, $2-\sqrt{x}$/1, $|x-1| - 4$/1, $f(x)$/1 }
	           {$0$, $\sqrt{3}$  , $4$, $5$   , $+\infty$}
	\tkzTabLine{, - , z , + , t , + , t , + }
	\tkzTabLine{, + , t , + , z , - , t , - }
	\tkzTabLine{, - , t , - , t , - , z , + }
	\tkzTabLine{, + , z , - , z , + , z , - }
\end{tikzpicture}


\subsubsection*{3.}
\begin{equation*}
	\begin{split}
		f(x) &= \ln(x+3) + \ln(x+2) - \ln(x+11) \\
		&= \ln \frac{(x+3)(x+2)}{x+11}
	\end{split}
\end{equation*}

\textbf{Domaine de définition}

$\ln$ étant défini sur $\mathbb{R}^{+}_{*}$ :

$D = ]-3 ; +\infty[ \cap ]-2 ; +\infty[ \cap ]-11 ; +\infty[ = ]-2 ; +\infty[$ (même si la forme rassemblée de $f$ est définie sur $]-11 ; -3[$)

\textbf{Valeurs et limites}

$\lim_{x \rightarrow -2^+} f(x) = -\infty$

$\lim_{x \rightarrow +\infty} f(x) = +\infty$

$f(x)=0 \Longleftrightarrow  \frac{(x+3)(x+2)}{x+11} = 1 \Longleftrightarrow x^2 +4x -5 =0 \Longleftrightarrow x \in [-5 ; 1] \Longleftrightarrow x=1$ sur $D$.

$f(0) = \ln\frac{6}{11} < 0$

\textbf{Tableau de signe}

\begin{tikzpicture}
	\tkzTabInit{$x$/1, $f(x)$/1 }
	           {$-2$, $1$ , $+\infty$}
	\tkzTabLine{d, - , 0 , + ,  }
\end{tikzpicture}


\subsection*{Exercice 2.17}

\subsubsection*{1.a.}
\begin{displaymath}
	\lim_{x \rightarrow 0^{+}} x e^{-\ln x} = \lim_{x \rightarrow 0^{+}} x x^{-1} = 1
\end{displaymath}

\subsubsection*{1.b.}
\begin{displaymath}
	\lim_{x \rightarrow 0} \frac{x+1}{x-1} = -1
\end{displaymath}

\subsubsection*{1.c.}
\begin{displaymath}
		\lim_{x \rightarrow 0} \frac{e^{\frac{1}{x^2}}+1}{2e^{\frac{1}{x^2}}-1}
		        = \lim_{x \rightarrow 0} \frac{1 + e^{\frac{-1}{x^2}}}{2-e^{\frac{-1}{x^2}}} = \frac{1}{2}
		        \quad \text{car } \lim_{x \rightarrow 0} e^{\frac{-1}{x^2}} = 0
\end{displaymath}

\subsubsection*{2.a.}
\begin{displaymath}
	\lim_{x \rightarrow +\infty} \frac{\sqrt{|x|}+1}{x-1}
		= \lim_{x \rightarrow +\infty} \frac{1 + \frac{1}{\sqrt{|x|}}}{\frac{x}{\sqrt{|x|}}-\frac{1}{\sqrt{|x|}}}
		= 0^{+}
\end{displaymath}

\begin{displaymath}
	\lim_{x \rightarrow -\infty} \frac{\sqrt{|x|}+1}{x-1} = 0^{-}
\end{displaymath}

\subsubsection*{2.b.}
\begin{displaymath}
	\lim_{x \rightarrow +\infty} \frac{x^3+1}{x^2-1} = \lim_{x \rightarrow +\infty} \frac{ x+ \frac{1}{x^2} }{ 1- \frac{1}{x^2}} =+\infty
\end{displaymath}

\begin{displaymath}
	\lim_{x \rightarrow -\infty} \frac{x^3+1}{x^2-1} = -\infty
\end{displaymath}

\subsubsection*{2.c.}
\begin{displaymath}
	\lim_{x \rightarrow +\infty} \frac{e^{2x}+1}{e^{3x}-1} = \lim_{x \rightarrow +\infty} \frac{1+ e^{-2x}}{e^{x}-e^{-2x}} = 0
\end{displaymath}

\begin{displaymath}
	\lim_{x \rightarrow -\infty} \frac{e^{2x}+1}{e^{3x}-1} = -1
\end{displaymath}


\subsubsection*{2.d.}
\begin{displaymath}
	\lim_{x \rightarrow +\infty} \left( \sqrt{|x|+1} - \sqrt{|x|-1} \right)
		= \lim_{x \rightarrow +\infty} \frac{|x|+1-(|x|-1)}{\sqrt{|x|+1} + \sqrt{|x|-1}} = \lim_{x \rightarrow +\infty} \frac{2}{\sqrt{|x|+1} + \sqrt{|x|-1}} = 0
\end{displaymath}
\begin{displaymath}
	\lim_{x \rightarrow -\infty} \left( \sqrt{|x|+1} - \sqrt{|x|-1} \right) = 0
\end{displaymath}

\subsubsection*{2.e.}
\begin{displaymath}
	\lim_{x \rightarrow +\infty} 2^{-x}e^x = \lim_{x \rightarrow +\infty} e^{x(1-\ln 2)} = +\infty \text{ cf. exercice 2.11 ou } e>2 \Longleftrightarrow 1 > \ln 2
\end{displaymath}

\begin{displaymath}
	\lim_{x \rightarrow -\infty} 2^{-x}e^x = \lim_{x \rightarrow -\infty} e^{x(1-\ln 2)} = 0^{}
\end{displaymath}


\subsection*{Exercice 2.18}

\subsubsection*{2.18.1.a}
\begin{displaymath}
	\lim_{x \rightarrow 0} e^{\frac{1}{x^2}} - e^{\frac{2}{x^2}}  = \lim_{x \rightarrow 0} e^{\frac{1}{x^2}}(1 - e^{\frac{1}{x^2}}) = -\infty
\end{displaymath}

\subsubsection*{2.18.1.b}
\begin{displaymath}
	\lim_{x \rightarrow 0} \frac{1}{x^2}  + \ln |x| = \lim_{x \rightarrow 0} \frac{1}{x^2}  (1+ x^2\ln |x|) = +\infty
\end{displaymath}
car $ \lim_{x \rightarrow 0} x^2\ln |x| = 0$ d'après la proposition 2.39

\subsubsection*{2.18.1.c}
\begin{displaymath}
	\lim_{x \rightarrow 0} \frac{1+e^{-x}}{(e^x-1)^2 - |x|\ln |x|} = + \infty
\end{displaymath}

car $ \lim_{x \rightarrow 0} (e^x-1)^2 = 0^{+}$ et $ \lim_{x \rightarrow 0} |x|\ln |x| = 0^{-}$, la limite du dénominateur est donc $0^{+}$

\subsubsection*{2.18.2.a}
\begin{displaymath}
	\lim_{x \rightarrow +\infty} \frac{\ln x + xe^{-x}-e^x}{e^x+1} = \lim_{x \rightarrow +\infty} \frac{-1 + e^{-x}\ln x + xe^{-2x}}{1 +e^{-x}} = -1
\end{displaymath}

d'après les propriétés de croissance comparée $\lim_{x \rightarrow +\infty} xe^{-2x} = \lim_{x \rightarrow +\infty} e^{-2x} = 0$ et 
$\lim_{x \rightarrow +\infty} e^{-x}\ln x = \lim_{x \rightarrow +\infty} e^{-x} =0 $.

\subsubsection*{2.18.2.b}
\begin{displaymath}
	\lim_{x \rightarrow +\infty} x^{\frac{1}{x}} = \lim_{x \rightarrow +\infty} e^\frac{\ln x}{x} = 1
\end{displaymath}

car $\lim_{x \rightarrow +\infty} \frac{\ln x}{x} = \lim_{x \rightarrow +\infty} \frac{1}{x} = 0$

\subsubsection*{2.18.2.c}
\begin{displaymath}
	\lim_{x \rightarrow +\infty} \frac{e^x - \ln x}{xe^x-x^2} = \lim_{x \rightarrow +\infty} \frac{1 - e^{-x}\ln x}{x-x^2e^{-x}} = 0
\end{displaymath}

\subsubsection*{2.18.2.d}
\begin{displaymath}
	\lim_{x \rightarrow +\infty} 1+ \frac{\ln (e^{3x})}{\sqrt{x}} = \lim_{x \rightarrow +\infty} 1+ \frac{3x}{\sqrt{x}} = +\infty
\end{displaymath}


\subsection*{Exercice 2.19}
\subsubsection*{2.19.1.a}
\begin{displaymath}
	f(x) = \frac{x-2}{x+1}
\end{displaymath}
\begin{equation*}
	\begin{split}
		f'(x) &=\frac{x+1 - (x-2)}{(x+1)^2}\\
		      &=\frac{3}{(x+1)^2}
	\end{split}	
\end{equation*}

\subsubsection*{2.19.1.b}
\begin{displaymath}
	f(x) = \frac{x^2-2x+1}{x^3+3x+3x+1} = \frac{(x-1)^2}{(x+1)^3}
\end{displaymath}
\begin{equation*}
	\begin{split}
		f'(x) &=\frac{2(x-1)(x+1)^3 - 3(x+1)^2(x-1)^2}{(x+1)^6}\\
		      &=\frac{2(x-1)(x+1) - 3(x-1)^2}{(x+1)^4} \\
		      &=\frac{(x-1)(2x+2-3x+3)}{(x+1)^4} \\
		      &=\frac{(x-1)(-x+5)}{(x+1)^4}
	\end{split}	
\end{equation*}

\subsubsection*{2.19.1.c}
\begin{displaymath}
	f(x) = x^{\frac{1}{3}} \ln x
\end{displaymath}
\begin{equation*}
	\begin{split}
		f'(x) &= \frac{1}{3} x^{-\frac{2}{3}} \ln x + x^{\frac{1}{3}} \frac{1}{x} \\
		      &=\left(1 + \frac{1}{3}\ln x \right) x^{-\frac{2}{3}}
	\end{split}	
\end{equation*}

\subsubsection*{2.19.1.d}
\begin{displaymath}
	f(x) = \frac{x}{e^x-2}
\end{displaymath}
\begin{equation*}
	\begin{split}
		f'(x) &= \frac{e^x-2 - xe^x}{(e^x-2)^2} \\
		      &= \frac{(1-x)e^x-2}{(e^x-2)^2} 
	\end{split}	
\end{equation*}

\subsubsection*{2.19.1.f}
\begin{displaymath}
	f(x) = \frac{e^x}{x^4 + 2x^2 +1} =\frac{e^x}{(x^2 +1)^2}
\end{displaymath}
\begin{equation*}
	\begin{split}
		f'(x) &= \frac{e^x(x^2+1)-2 \times 2xe^x}{(x^2 +1)^3} \\
		      &= \frac{x^2 - 4x + 1}{(x^2 +1)^3}e^x
	\end{split}	
\end{equation*}



\subsubsection*{2.19.2.a}
\begin{displaymath}
	f(x) = \ln^2 x
\end{displaymath}
\begin{equation*}
	\begin{split}
		f'(x) &= 2\frac{\ln x}{x}
	\end{split}	
\end{equation*}

\subsubsection*{2.19.2.b}
\begin{displaymath}
	f(x) = \frac{\sqrt{x+1}}{x^2-1} = \frac{\sqrt{x+1}}{(x-1)(x+1)}= \frac{1}{\sqrt{x+1}(x-1)}
\end{displaymath}
\begin{equation*}
	\begin{split}
		f'(x) &= -\frac{\frac{x-1}{2\sqrt{x+1}} + \sqrt{x+1}}{(x+1)(x-1)^2} \\
		      & = -\frac{3x+1}{2(x+1)^{\frac{3}{2}} (x-1)^2}
	\end{split}	
\end{equation*}

\subsubsection*{2.19.2.c}
\begin{displaymath}
	f(x) = \sqrt{e^{x^3}}
\end{displaymath}
\begin{equation*}
	\begin{split}
		f'(x) &= \frac{3x^2 e^{x^3}}{2\sqrt{e^{x^3}}} \\
		      &= \frac{3}{2}x^2 \sqrt{e^{x^3}}
	\end{split}	
\end{equation*}


\subsection*{Exercice 2.20}
\begin{displaymath}
	\lim_{x \rightarrow 0^{+}} (1+x)^\frac{1}{x}
		= \lim_{x \rightarrow 0^{+}} e^{\frac{\ln(1+x)}{x}}
		= \lim_{x \rightarrow 0^{+}} e^{1} 
		= e
\end{displaymath}

\begin{displaymath}
	\lim_{x \rightarrow 0^{+}} (1+x)^\frac{1}{x^2} 
		= \lim_{x \rightarrow 0^{+}} e^{\frac{1}{x}\frac{\ln(1+x)}{x}} 
		= \lim_{x \rightarrow 0^{+}} e^{\frac{1}{x}}
		= +\infty
\end{displaymath}

\begin{displaymath}
	\lim_{x \rightarrow 0^{+}} (1+x)^\frac{1}{\sqrt{x}} 
		= \lim_{x \rightarrow 0^{+}} e^{\frac{\ln(1+x)}{\sqrt{x}}} 
		= \lim_{x \rightarrow 0^{+}} e^{\sqrt{x}\frac{\ln(1+x)}{x}}
		= \lim_{x \rightarrow 0^{+}} e^{\sqrt{x}}
		= 1
\end{displaymath}

car $\lim_{x \rightarrow 0^{+}} \frac{\ln(1+x)}{x} = 1$ (cf. dérivée de $\ln$ en $1$)

\subsection*{Exercice 2.21}

$f(x) = x^x = e^{x \ln x}$

Domaine de définition : $D = \mathbb{R}^+_*$

\begin{displaymath}
	\lim_{x \rightarrow 0^{+}} f(x) = 1  \text{ (cf. croissance comparée de $x$ et $\ln x$)}
\end{displaymath}

\begin{displaymath}
	\lim_{x \rightarrow +\infty} f(x) = +\infty
\end{displaymath}

\begin{displaymath}
	f'(x) = (1+ \ln x)e^{x \ln x} = (1+ \ln x)x^x
\end{displaymath}

\begin{displaymath}
	f'(x) = 0 \Longleftrightarrow 1 + \ln x = 0 \Longleftrightarrow x = \frac{1}{e} \approx 0,368
\end{displaymath}

\begin{displaymath}
	f\left(\frac{1}{e}\right) = e^{-\frac{1}{e}} \approx 0,692
\end{displaymath}

\begin{tikzpicture}
	\tkzTabInit{$x$/1, $f'$/1, $f$/2 }
	{$0$  , $\frac{1}{e}$   , $+\infty$}
	\tkzTabLine{d, - , z , +, }
	\tkzTabVar{+C/$1$, -/$e^{-\frac{1}{e}}$ , +/$\infty$ }
\end{tikzpicture}

\begin{tikzpicture}
	\begin{axis}[
		axis lines=center,
		height=5cm,
		width=10cm,
		xlabel={$x$},
		ylabel={$f(x)$},
		domain=0:1.1,
		ymin=0.5,
		]
		\addplot[blue, samples=101]{x^x};
	\end{axis}	
\end{tikzpicture}

\subsection*{Exercice 2.22}

\subsubsection*{Exercice 2.22.1}

$f_1(x) = \frac{1}{1+e^{-x}}$

\begin{displaymath}
	\lim_{x \rightarrow +\infty} f_1(x) = 1
\end{displaymath}
\begin{displaymath}
	\lim_{x \rightarrow -\infty} f_1(x) = 0
\end{displaymath}

$e^{-x}$ décroissante donc $(1+e^{-x})$ décroissante, donc $f_1$ est croissante.

$f_1(0) = \frac{1}{2}$

\begin{displaymath}
	f_1'(x) = \frac{e^{-x}}{(1+e^{-x})^2} > 0
\end{displaymath}

Ce qui confirme que $f_1$ est strictement croissante sur $\mathbb{R}$.

\begin{displaymath}
	f_1''(x) = \frac{-(1+e^{-x})+2}{(1+e^{-x})^3}e^{-x}
	         = \frac{1-e^{-x}}{(1+e^{-x})^3}e^{-x}
\end{displaymath}

$f_1$ présente donc un point d'inflexion en $0$.

On constate que $f_1(x) + f_1(-x) = \frac{1}{2}$ ce qui signifie que la représentation
graphique est symétrique par rapport au point $(0 ; \frac{1}{2})$


\begin{tikzpicture}
	\tkzTabInit{$x$/1, $f'$/1, $f$/2 }
	{$-\infty$  , $+\infty$}
	\tkzTabLine{, + , }
	\tkzTabVar{-/$0^{+}$, +/$1^-$ }
	\tkzTabVal{1}{2}{0.5}{0}{$\frac{1}{2}$}
\end{tikzpicture}

\begin{tikzpicture}
	\begin{axis}[
		axis lines=center,
		height=5cm,
		width=10cm,
		xlabel={$x$},
		ylabel={$f(x)$},
		domain=-5:5,
		ymin=0,
		]
		\addplot[blue, samples=101]{1/(1+exp(-x))};
	\end{axis}	
\end{tikzpicture}


\subsubsection*{Exercice 2.22.2}

$f_2(x) = 2x\sqrt{1-x^2}$ sur $[-1 ; 1]$

$f_2$ est impaire

\begin{displaymath}
	f_2'(x) = 2\sqrt{1-x^2} + \frac{2x(-2x)}{2\sqrt{1-x^2}} 
	        = \frac{2(1-x^2)-2x^2}{\sqrt{1-x^2}}
	        = \frac{2-4x^2}{\sqrt{1-x^2}}
\end{displaymath}

$f_2$ s'annule donc en $\pm \frac{1}{\sqrt{2}}$

$f\left( \frac{1}{\sqrt{2}} \right) = 1$
	
\begin{tikzpicture}
	\tkzTabInit{$x$/1, $f'$/1, $f$/2 }
	{$-1$, $-\frac{1}{\sqrt{2}}$, $\frac{1}{\sqrt{2}}$  , $1$}
	\tkzTabLine{, + , z , -, z, + }
	\tkzTabVar{+/$0$, -/$-1$,+/$1$, -/$0$ }
\end{tikzpicture}

\begin{tikzpicture}
	\begin{axis}[
		axis lines=center,
		height=5cm,
		width=10cm,
		xlabel={$x$},
		ylabel={$f(x)$},
		domain=-1:1,
		ymin=-1,
		]
		\addplot[blue, samples=101]{2*x*sqrt(1-x*x)};
	\end{axis}	
\end{tikzpicture}

\subsubsection*{Exercice 2.22.3}

$f_3(x) = \exp\left( \frac{x-1}{x+1}\right) = \exp\left(1- \frac{2}{x+1}\right)$

$D = \mathbb{R} \setminus \{-1\}$

$f_3$ est strictement positive sur $D$

\begin{displaymath}
	\lim_{x \rightarrow -\infty} f_3(x) = e^{+}
\end{displaymath}

\begin{displaymath}
	\lim_{x \rightarrow -1^{-}} f_3(x) = +\infty
\end{displaymath}

\begin{displaymath}
	\lim_{x \rightarrow -1^{+}} f_3(x) = 0
\end{displaymath}

\begin{displaymath}
	\lim_{x \rightarrow +\infty} f_1(x) = e^{-}
\end{displaymath}

\begin{displaymath}
	f_3'(x) = \frac{x+1-(x-1)}{(x+1)^2} \exp\left( \frac{x-1}{x+1}\right)
	        = \frac{2}{(x+1)^2} \exp\left( \frac{x-1}{x+1}\right) > 0
\end{displaymath}

\begin{displaymath}
	f_3''(x) = \left(\frac{-4(x+1)}{(x+1)^4} + \frac{4}{(x+1)^4}\right) \exp\left( \frac{x-1}{x+1}\right)
	         = \frac{-4x}{(x+1)^4} \exp\left( \frac{x-1}{x+1}\right)
\end{displaymath}

$f_3$ présente un point d'inflexion en $(0 ; \frac{1}{e})$

\begin{tikzpicture}
	\tkzTabInit{$x$/1, $f'$/1, $f$/2 }
	{$-\infty$, $-1$, $+\infty$}
	\tkzTabLine{, + , d , + }
	\tkzTabVar{-/$e$, +DC-/$+\infty$/$0$, +/$e$ }
\end{tikzpicture}



\begin{tikzpicture}
	\begin{axis}[
		axis lines=center,
		height=10cm,
		width=10cm,
		xlabel={$x$},
		ylabel={$f(x)$},
		ymin=0,
		]
		\addplot[blue, domain=-10:-2.5, samples=51]{exp((x-1)/(x+1))};
		\addplot[blue, domain=-1:9, samples=51]{exp((x-1)/(x+1))};
		\addplot[red, thin, domain=-10:9]{exp(1)};
		\addplot[red] coordinates {(-1,0)(-1,10)};
	\end{axis}	
\end{tikzpicture}

\end{document}]